\section{Glucose và saccharose}

\subsection{Carbohydrate}

Carbohydrate (cacbohiđrat) là nhóm hợp chất hữu cơ có thành phần nguyên tố chủ yếu gồm carbon, hydrogen và oxygen. Nhiều carbohydrate có thể biểu diễn bằng công thức dạng \(\ce{C_n(H2O)_m}\), nhưng đây chỉ là cách biểu diễn khái quát, không phải công thức cấu tạo chung cho mọi chất.

Glucose, saccharose, tinh bột và cellulose là những carbohydrate quen thuộc trong đời sống. Chúng có vai trò quan trọng trong dinh dưỡng và là nguyên liệu cho nhiều quá trình sản xuất.

\subsection{Glucose}

Glucose có công thức phân tử \(\ce{C6H12O6}\). Glucose có trong quả chín, mật ong và một số bộ phận của thực vật. Đây là chất rắn kết tinh, không màu, có vị ngọt và tan tốt trong nước.

Glucose là chất dinh dưỡng quan trọng đối với người và động vật. Dung dịch glucose được sử dụng trong thực phẩm và trong một số chế phẩm y tế; chẳng hạn dung dịch glucose 5\% có 5 gam glucose trong 100 mL dung dịch.

\subsubsection*{Phản ứng tráng bạc}

Khi đun nóng glucose với thuốc thử Tollens (dung dịch \(\ce{AgNO3}\) trong ammonia), ion bạc được khử thành bạc kim loại. Bạc tạo thành bám lên thành ống nghiệm, tạo lớp sáng như gương.

Có thể biểu diễn phản ứng bằng phương trình phân tử giản lược:
\[
\ce{C6H12O6 + Ag2O ->[NH3,\ t^\circ] C6H12O7 + 2Ag}.
\]

Theo phương trình, 1 mol glucose tạo ra 2 mol bạc. Tỉ lệ này là cơ sở để giải các bài toán về tráng gương, tráng ruột phích hoặc tạo lớp bạc trên bề mặt vật liệu.

\subsubsection*{Phản ứng lên men rượu}

Khi có enzyme của men rượu và điều kiện thích hợp, glucose lên men tạo ethylic alcohol và carbon dioxide:
\[
\ce{C6H12O6 ->[enzyme] 2C2H5OH + 2CO2}.
\]

Phản ứng này là một mắt xích trong quá trình sản xuất đồ uống lên men và nhiên liệu sinh học. Khi giải bài toán, cần chú ý tỉ lệ mol giữa glucose, ethylic alcohol và carbon dioxide theo phương trình.

\begin{ghinho}{Cần nhớ về glucose}
\begin{itemize}
\item Công thức phân tử: \(\ce{C6H12O6}\).
\item Có trong quả chín, mật ong; tan tốt trong nước và có vị ngọt.
\item Có phản ứng tráng bạc và phản ứng lên men rượu.
\item Trong phản ứng tráng bạc, 1 mol glucose tạo 2 mol \ce{Ag}.
\end{itemize}
\end{ghinho}

\subsection{Saccharose}

Saccharose có công thức phân tử \(\ce{C12H22O11}\), có nhiều trong mía, củ cải đường và một số loại quả. Saccharose là chất rắn kết tinh, không màu, có vị ngọt và tan tốt trong nước. Đây là nguyên liệu quan trọng trong công nghiệp thực phẩm.

\subsubsection*{Phản ứng thủy phân}

Khi đun nóng saccharose trong môi trường acid hoặc có enzyme thích hợp, saccharose bị thủy phân tạo glucose và fructose:
\[
\ce{C12H22O11 + H2O ->[H^+,\ t^\circ] C6H12O6 + C6H12O6}.
\]

Glucose và fructose tạo thành có cùng công thức phân tử \(\ce{C6H12O6}\), nhưng là hai chất khác nhau. Không được hiểu phương trình trên là saccharose tạo ra hai phân tử glucose.

Saccharose không cho phản ứng tráng bạc khi thử trực tiếp trong điều kiện thông thường. Sau khi thủy phân, hỗn hợp sản phẩm có thể tham gia phản ứng tráng bạc.

\subsection{Phân biệt glucose và saccharose}

Nếu hai mẫu dung dịch không màu lần lượt chứa glucose và saccharose, có thể dùng thuốc thử Tollens để phân biệt:
\begin{itemize}
\item Mẫu tạo lớp bạc bám lên thành ống nghiệm khi đun nóng là mẫu chứa glucose.
\item Mẫu không xuất hiện lớp bạc trong phép thử trực tiếp là mẫu chứa saccharose.
\end{itemize}

Không thể phân biệt hai mẫu chỉ bằng vị ngọt hoặc màu sắc vì cả hai đều là chất rắn kết tinh không màu, có vị ngọt và tan tốt trong nước.

\begin{vidu}
Một mẫu dung dịch không màu là glucose, mẫu còn lại là saccharose. Hãy nêu cách nhận biết hai mẫu bằng phương pháp hóa học và viết phương trình phản ứng xảy ra.
\end{vidu}
\begin{loigiai}
Lấy một lượng nhỏ mỗi mẫu vào hai ống nghiệm riêng. Thêm thuốc thử Tollens vào từng ống nghiệm rồi đun nóng nhẹ trong điều kiện phù hợp.

Ống nghiệm xuất hiện lớp bạc bám trên thành là mẫu chứa glucose; mẫu còn lại là saccharose.

Phương trình phản ứng:
\[
\ce{C6H12O6 + Ag2O ->[NH3,\ t^\circ] C6H12O7 + 2Ag}.
\]

\textbf{Vậy,} có thể nhận biết glucose bằng phản ứng tráng bạc; saccharose không cho hiện tượng này khi thử trực tiếp.
\end{loigiai}
